% Experimental setup. Sources, in order of authority: docs/research_findings.md F108-F151, AGENT/handoff.md; the code
% (configs/space_runtime.yaml, openral_ext/capability_server/{contracts.yaml,node.py}, experiments/operator_console/
% console.py, space_context/{link.py,earth_visibility.py,profiles/earth_link.json}, evaluation/{interaction,violations,
% rows,tables}.py); docs/heldout_design.md §1, §2, §6, §8, §9, §14; docs/preregistration.md;
% docs/baselines/authoring_spec.md; docs/reauthoring/spec.md; docs/exp2/extension_spec.md; F148 for the ablations.
% Every number carries its source in a comment next to it.
\providecommand{\TBD}[1]{\textcolor{red}{[TBD: #1]}}

\section{Experimental setup}
\label{sec:setup}

\subsection{Simulated lunar base and robot}
\label{sec:setup-site}

The experiments run in a simulated lunar base on a LOLA elevation model of a ridge near the south pole, at
89.455\,$^\circ$S, 213.64\,$^\circ$E \citep{barker2021lola}. The simulator is Isaac Sim 5.1 with the Space Robotics Bench
\citep{orsula2024srb} (Fig.~\ref{fig:site}). Five serviced units each have their own load and criticality. The power
distribution units PDU-1, PDU-2 and PDU-3 feed habitat life support (crew-critical), the communication station
(mission-critical) and the science instruments (non-critical). The battery storage unit (BESS) is mission-critical, and
the ISRU oxygen plant is non-critical. Each unit has a control panel with a reset button and a status lamp, and a bay
that holds a replaceable module. The PDUs and the BESS also have a radiator or condenser behind a shutter with its own
button. Stands at the depot, the landing pad, the habitat and the return area hold spares, logistics items and returns.
% 89.455 S 213.64 E: F108 §5 (earth_visibility.py SITE: -89.4553, 213.642). Isaac Sim 5.1: docs/upstream_pins.md.
% Units, loads, criticality: interaction.json `units`; docs/heldout_design.md §9 "As built".

One site model governs every scenario, and scenarios differ only in their initial state. A press on a unit's panel
resets its breaker. A drawn shutter makes the unit trip on over-temperature 30\,s after a reset. An unserviceable module
makes it trip after a 3\,s self-test, and an empty bay makes it trip at once. Each unit reports its own telemetry: breaker
state, trip cause, energized state, load, criticality, reset count and the serial of each installed module. Mission control's inventory record lists every tracked item by serial, recorded
location and status. Assets are rigid and actuators are push buttons. The carried modules are look-alikes, far lighter
than flight hardware.
% 30 s: F112 §4(d) retrip_delay_s "from a stated thermal assumption"; 3 s: F108 §1, site.py SELFTEST_S = 3.0.
% "far lighter": docs/heldout_design.md §9 honesty note (the 1 kg look-alikes).

The robot is LunarBot, a rover with four independently steered wheels. It carries a seven-joint RM-75 arm, an EG2
parallel gripper, a Livox Mid-360 lidar, and RGB-D cameras on its mast and wrist. Its software runs on ROS~2 Humble
\citep{macenski2022ros2}. It localizes with AMCL on a map surveyed with its own lidar and drives with Nav2
\citep{macenski2020nav2}. The arm plans collision-checked motions with MoveIt \citep{coleman2014moveit}. Pick, place and
press are OpenRAL skills \citep{openral}, and pick and press ask the shared vision-language model which surface in the
wrist view is the target. All methods call the skills through one capability boundary with the same names and
arguments, and the boundary returns only robot-observable evidence. Every module carries a Data Matrix data plate
\citep{nasa2020std6002}, which only the scanner added in Experiment II reads.
% Robot: docs/lunar_task_catalog.md header; docs/lunar_bot_capability_set_plan.md (Mid-360); contracts.yaml (D435i
% RGB-D, RM-75 joints). Boundary: openral_ext/capability_server/node.py docstring. Plates: heldout_design §4, §6.

The link to Earth is Direct-with-Earth only. Earth's elevation at the site comes from SPICE and the DE440 ephemeris
\citep{acton1996spice,park2021de440}, over a flat horizon. The mission epoch, 2028-10-11 01:18:39 UTC, puts Earth-set
925\,s of mission time after the start of every run. That interval is three request-and-answer cycles of the protocol
in Section~\ref{sec:setup-operator}, fixed before any run of a case that needs the ground. Earth rises again 13.4 days
later. The one-way delay is 4.27\,s, the light time of 1.27\,s plus 3\,s of ground latency at the low end of the lunar
range \citep{landon2025delay}. Mission time is the simulator clock. The simulation ran at about one fifth of real time,
and the reasoning model's latency passed into mission time at the same ratio. Mission time therefore charges the model's
deliberation at about one fifth of its wall-clock duration.
% Epoch, 925 s, 13.4 d (aos_after_start_s 1161258.8 s), 4.27 = 1.271 + 3.0 s: space_context/profiles/earth_link.json,
% F116. "three request-and-answer cycles": 3 x (4.27 + 300 + 4.27) s, earth_visibility.py MISSION_START_BEFORE_LOS_S.
% "about one fifth": F116 "deliberation runs at the sim/wall ratio (~0.2)"; handoff "sim/wall ~0.2 with six at once".
% docs/preregistration.md §3 asks for each request's wall-clock send time beside its mission time (open question).

\begin{figure}[t]
  \centering
  \fbox{\parbox[c][4cm][c]{0.9\columnwidth}{\centering site figure placeholder}}
  \caption{\TBD{The simulated lunar base: the five serviced units, the stands, the dock, and LunarBot at a unit's
  bay. Third-person render from the scene.}}
  \label{fig:site}
\end{figure}

\subsection{Missions and the novelty ladder}
\label{sec:setup-missions}

Two mission families share the site. Mission A is the restoration objective quoted in Section~\ref{sec:requirements},
with the reported unit as a parameter. Mission B reads ``Deliver the \emph{item} staged at the \emph{source} (slot 0
per the manifest) to \emph{destination}, then return the rover to the dock with the arm stowed.'' The external evaluator
scores a Mission A run as successful when the faulted unit is restored and still restored at the end, with the rover
docked and the arm stowed. L2-A also needs the shutter open. L3 and L4 need a serviceable module in the target bay and
the failed one out of it. A Mission B run succeeds when the named item rests on the destination inside its footprint,
with the rover docked and the arm stowed.
% Mission B brief: lunar_base_mission_b/task.py mission_brief(). Success conditions: docs/heldout_design.md §4.

\begin{table*}[t]
  \centering
  \caption{The thirteen instances of Experiment I, by novelty level. L0 and L1 are covered for both baselines; L2-A to
  L4 are held-out. Each held-out situation has a documented precedent.}
  \label{tab:instances}
  \small
  \begin{tabular}{lllp{0.36\textwidth}l}
    \toprule
    level & frozen baseline lacks & mission & instance & precedent \\
    \midrule
    L0 & nothing & A & PDU-2 trips; reset & \citep{lewis2022uncrewed} \\
       &         & B & spare ORU, depot stand to the PDU-3 spares shelf & \citep{lewis2022uncrewed} \\
    L1 & variable values & A & BESS trips; reset & \citep{lewis2022uncrewed} \\
       &                 & A & ISRU plant trips; reset & \citep{lewis2022uncrewed} \\
       &                 & A & the report names PDU-2; the BESS is the tripped unit & \citep{nasa2015issdsr} \\
       &                 & B & spare battery module, depot spares stand to the BESS receiving stand & \citep{lewis2022uncrewed} \\
       &                 & B & consumables container, landing-pad stand to the habitat logistics stand & \citep{lewis2022uncrewed} \\
    L2-A & a cause & A & PDU-1: radiator shutter drawn, over-temperature after each reset & \citep{gaier2005apollo,nasa1971apollo15} \\
         &         & A & BESS: condenser shutter drawn, over-temperature after each reset & \citep{gaier2005apollo,nasa1971apollo15} \\
    L2-B & a relation & B & habitat logistics stand occupied by an outbound return container & \citep{nasa2019issdsr0617} \\
    L3 & a composition & A & BESS battery module fails its self-test; spare on the depot spares stand & \citep{dextre2011rpcm,nasa2019issdsr0402} \\
       &               & A & ISRU filter cartridge clogged; spare on the ISRU receiving stand & \citep{ogs2011filters} \\
    L4 & a goal conflict & A & PDU-1 module failed, the depot spare recorded failed; the only serviceable module is in energized PDU-3 & \citep{cfr1214403,gao1988cannibalization} \\
    \bottomrule
  \end{tabular}
\end{table*}
% Instances: docs/heldout_design.md §14 (13 instances, 7 covered + 6 held-out). Precedents: §2, §9, §13.

The novelty ladder grades the held-out cases by what the frozen baseline model lacks (Table~\ref{tab:instances}). L0
and L1 differ from a covered procedure in nothing or only in variable values, so they are covered for both baselines.
At L2-A a drawn shutter makes every reset trip again, a cause the baselines do not model. At L2-B the destination is
already occupied, a relation they do not model. L3 needs a composition across the two families. Both baselines contain
every piece of its solution, the reset of Mission A and the pick, carry and place of Mission B, but not their
combination. L4 poses a goal conflict. The only serviceable module sits in an energized, non-critical unit, and taking it
needs mission control's authorization before Earth sets. Each level runs on one or two units. Two precedents hold only
for the category, and the drawn shutter and the occupied stand are our extrapolations from them.
% docs/heldout_design.md §2 (ladder, L3 "the sharpest", L4 authority + link), §14 ("each level is instantiated on one
% or two units"); §2 and §9 honesty notes on the two extrapolations.

Physical actions take only forms verified before the campaign: pressing a button, picking and placing a top-handled item
on a horizontal surveyed support, driving and docking. Every held-out physical sequence ran once before the campaign as a
fixed chain of calls through the shared capabilities, with no method in the loop, and completed. A physical failure in
the campaign therefore cannot pass for a reasoning result.
% docs/heldout_design.md §1 rule 3 and §12; F112 §6-7 ("all nine held-out sequences now executed end to end").

\subsection{Compared methods}
\label{sec:setup-baselines}

The proposed runtime uses Qwen3.8-Max as its reasoning model, with thinking enabled and a 600\,s timeout per step. The
vision-language model is DeepSeek-V4.1-Flash, with a 300\,s timeout. It answers the proposed runtime's observation
queries and makes the target selection inside pick and press for all three methods, so it is the only model the
baselines use. No request sets a sampling parameter. The mission step budget is 80 steps, and the loop suppresses a
fourth identical call. The palette holds 27 tools in Experiment I. The guard has one target rule, the removal of a module
from an energized unit's bay, and a latched emergency stop prohibits every effect. No other authority class or hard
invariant is configured. Tables and figures label the proposed runtime Ours.
% configs/space_runtime.yaml: model Qwen3.8-Max, thinking true, timeout_s 600, retry_cap_per_call 3, max_steps 80,
% guard target_authority energized_equipment_removal, invariants [], fdir_prohibits estop; every contract routine.
% VLM: node.py VLM_MODEL DeepSeek-V4.1-Flash; 300 s: F114, F137. No sampling parameter: F137. 27 tools: F124.
% Step budget: Experiment I counted every proposal; Experiment II's final system counts only calls on the robot (F147).

The reasoning model was chosen among four candidates before the freeze, on development runs and offline replays of
recorded decisions. The proposed runtime was developed on development runs of the held-out instances under one rule. A
change must name no situation, place or serial, hold for every mission, and trace to a recorded finding. Before the
campaign, an audit checked every change against this rule and searched every text the model reads for the words of the
held-out cases. No development run is scored.
% F119-F123 (candidates DeepSeek-V4.1-Flash, Kimi-K3, Qwen3.8-Max, GLM-5.3-Flash), F137 and F139 (freeze audit),
% AGENT/todo.md "Ours: generic changes only, three tests". Open question: level of detail on the model choice.

PLEXIL is the official executive, Release 6.1.3 \citep{estlin2006plexil}. An interface adapter exposes each capability
as a Command and each query as a Lookup, and logs the issuing plan node with every call. EUROPA\,$\rightarrow$\,PLEXIL is
our integration of two NASA open-source components, not a NASA product. The EUROPA planner \citep{bedraxweiss2004europa}
holds one domain model for both missions and takes the units' telemetry, the site survey and the inventory record as
facts. A bridge dispatches its activities through the PLEXIL executive, and the planner replans from the updated facts
after each result. A one-line patch fixed a use-after-free in EUROPA's temporal network that crashed the planner on some
inputs. Another solver defect made the host solve one goal at a time, which makes the planner greedier, not stronger.
% PLEXIL 6.1.3: F133. Adapter node id: F133. EUROPA patch: F112 §1 ("10 memcheck errors -> 0"). One goal per solve:
% docs/baselines/authoring_spec.md change table ("makes EUROPA greedier, i.e. weaker").

Both baselines follow a specification locked before any further baseline run, so that no failed run could pull a
baseline stronger or weaker after the fact. A baseline is its native mechanism plus the knowledge of the covered
specification, authored natively for every covered instance. Mission A's formal goal is every unit restored, and the
plan or the solver finds the tripped unit from telemetry. Mission B is parameterized by item, source and destination.
Failure handling has three tiers. A failed call is repeated up to 3 attempts, the cap the proposed runtime has on an
identical call. A failed goal check that the covered brief asks for, such as a breaker that must still read closed at
the dock, re-runs its action under the same budget. The baseline then asks mission control and holds. No branch may key
on a cause or a situation, such as a trip cause, an inventory status, an occupied support or the stage at which a skill
failed. That knowledge is what the re-authoring study and Experiment II price. The operability standard for space
segments treats the retry count as a parameter and sends what a level cannot handle to the ground
\citep{ecss2008operability}.
% docs/baselines/authoring_spec.md "Locked specification" (lock commit 1b6a40f); retry = configs retry_cap_per_call 3;
% ECSS-E-ST-70-11C §5.7.5 ("no standard value"). Change control after the lock: shared-infrastructure and conformance
% fixes only.

A larger retry budget could have changed where the baselines stop. After the freeze, both baselines ran offline against
stand-ins of every held-out state with one, two, three and five attempts per call, 48 runs in all. PLEXIL's plans unroll
the budget, so each other budget ran on a derived copy of the locked plans. The copy for three attempts reproduces the
locked plans character for character.
% F136, F143 (48 offline runs: 2 baselines x 6 held-out instances x N = 1/2/3/5; derived(3) identical to the lock).

\subsection{Mission control}
\label{sec:setup-operator}

Mission control is a scripted protocol, published before the runs and identical for all methods. It answers a request
by its declared kind and by the calls the guard holds, never by its text. A flight-rule table decides an authorization.
It grants the removal of a module from a unit with a non-critical load while a crew-critical unit is down, and denies
anything else with a reason code \citep{cfr1214403}. An authority check receives the published rules. An information
request receives the statement that mission control holds no data beyond the robot's own queries. A new-procedure
request receives the instruction to hold in a safe state while a procedure is written. Historical new procedures took
from hours to days \citep{nasa2019mbsu3}, so nothing is uplinked during a run, and this answer alone marks a run as
requiring ground re-authoring. A request without a kind counts as a new-procedure request, and every baseline request
is of this kind.
% experiments/operator_console/console.py decide(), module docstring; F125; docs/heldout_design.md §8 table
% (historical times: Apollo 13 ~34 h, STS-120 4 d, MBSU-3 2019 3 d, Dextre RPCM ~13 months).

Every request is decided 300\,s of mission time after it reaches the ground. We chose this value from minute-scale
decisions in rover operations \citep{lim2024viper}, fixed it before the campaign and ran no sensitivity study. The human
intervention ratio scales linearly with it. While Earth is below the horizon, the robot's link node reports a request as
queued for the next pass. A queued new-procedure request still counts as re-authoring, and an authorization held back
past the end of the run counts as unanswered. Mission control's answer to an authorization is recorded at the ALFUS level
of a tactical goal, and a new procedure at the level of a mission task \citep{huang2007alfus}. No live operator took part
in a scored run, which keeps the ground identical across methods \citep{riek2012woz}.
% 300 s: console.py DECISION_TIME_S; docs/heldout_design.md §8 ("value is ours", "no sensitivity runs", HIR scales
% linearly). Queued escalation = re-authoring, unanswered authorization: docs/preregistration.md §3, F116.
% ALFUS labels: console.py decide() (tactical_goal / mission_task).

\subsection{Experiment I: autonomous mission generalization}
\label{sec:setup-exp1}

Experiment I runs the three frozen methods on the thirteen instances of Table~\ref{tab:instances}. The baselines' plans
and domain fix their decisions, so each baseline ran once per instance. The proposed runtime ran once per covered
instance and three times per held-out instance, to expose the variance of its reasoning model. That gives 51 scored runs.
The three methods of an instance ran on the same hardware, one simulation per GPU on a six-GPU machine. Every run
recorded its full raw data, so a later change of a metric's definition is recomputed, never re-simulated. Defects of the
shared layer found during the campaign were fixed for every method, and Experiment I was then repeated in full. Two runs
of the repeat, invalidated by a further shared defect, were rerun. The results report the repeat.
% Runs: docs/heldout_design.md §14 (26 baseline + 7 + 18 Ours = 51); F142 (51 scored runs, 6 GPUs; two runs
% invalidated by shared-layer defects, fixed and rerun). Raw data: docs/preregistration.md preamble. First full campaign:
% F140 (exp1), shared-layer fixes F141, repeat F142 (exp1b). Open question: whether to report exp1 in an appendix.

\subsection{Experiment II: extending the mission envelope}
\label{sec:setup-exp2}

Experiment II adds a capability after the freeze and prices the change. The capability, \texttt{scan\_label}, reads the
serial plates of the item on a surveyed support with the wrist camera. It was built before the experiment, supplied
identically to all methods and registered with none. Its specification says that the inventory record can be wrong and
that items of one kind differ only by serial. A Mission B brief can now name the cargo by serial, and the delivered unit
must match it. One engineer extended each method, a fresh Claude session \TBD{model and version} that saw only the
specification and a visible validation suite. The sessions ran in the order proposed runtime, PLEXIL,
EUROPA\,$\rightarrow$\,PLEXIL, each from the previous merged state. The baselines thus started with the first engineer's
capability contract, and EUROPA\,$\rightarrow$\,PLEXIL also with PLEXIL's adapter Lookups. This order favours the
baselines.
% docs/exp2/extension_spec.md §1-§4; F131 (built, unregistered), F144 (supplied before any timer), F146 (protocol as
% run: fresh Claude session per method, order Ours -> PLEXIL -> EUROPA from the merged state).

Each engineer changed only its method's own artifacts. A change to the proposed runtime had to name no situation, place
or serial and hold for every task. The baselines could encode scanning and searching, since the extension specification
joined their covered specification, but not a particular wrong record. A timer ran over active work, including waiting
on the engineer's own validation run, and stopped at the first pass of the suite. The suite is one live run of a visible
case with correct records. It checks the wiring, the evaluator's success, a depot scan that read the named serial, that
nothing was dropped and that the trace is complete.
% docs/exp2/extension_spec.md §4-§5; experiments/exp2/g3_validate.py (F144 "Visible validation suite (G3)").

Two hidden cases stayed outside the repository until the last method was tagged. In H1 the named spare stands on
PDU-1's spares shelf, and a failed unit of the same kind stands where the record puts the spare. An ISS cargo bag once
stayed lost for about a year while the inventory records pointed to the wrong modules \citep{berger2021rfid}. In H2 the
two depot units carry each other's records, as when a THEMIS flight model received a part of the right part number and
the wrong serial \citep{themis2007pfr196}. Each method ran each case once before its extension. After it, the proposed
runtime ran each case three times and each baseline once. These runs exposed three generic defects. Two were shared,
a scan view that could cut a plate at the frame's edge and a scan result that could not tell an empty support from an
unreadable plate. The third was the proposed runtime's step budget, which counted its own plan bookkeeping. The fixes
went into the extended systems, and only the two runs they decided were rerun.
% experiments/exp2/heldout/README.md (H1, H2, precedents.md); artifacts/eval/exp2/tables.md §2 (before 0/1 per case
% per method; after Ours 3/3 per case, baselines 1/1). F145 (scan view, rerun EUROPA H1), F147 (on-support fact,
% step budget, rerun Ours H1 r2).

\subsection{Re-authoring cost of the held-out cases}
\label{sec:setup-reauthoring}

Every held-out baseline run in Experiment I ended with a request for a new procedure. The re-authoring study measures
that procedure. For each pair of baseline and held-out case, eight in all, a fresh Claude session \TBD{model and version}
started from the Experiment I freeze. It received the diagnosis that mission control's engineers would have after the
escalation, meaning the cause, the facts that matter and the goal, but not the procedure. Tier three of the lock was
lifted for that case only, so a branch or a domain relation could key on its cause. The retry budget and the covered
procedure stayed. The engineer could not open any record of the proposed runtime. A timer with a cap of 3\,h ran over
active work and stops at the first passing validation. Validation runs the modified baseline live on every instance of
the case. It passes when the evaluator scores every instance a success, nothing was dropped and no module left an
energized unit without approval. The proposed runtime needed no re-authoring for these cases, so its cost here is zero.
% docs/reauthoring/spec.md; F149 (eight sessions, opus; briefs docs/reauthoring/cases/{l2a,l2b,l3,l4}.md); F151
% (authoring done offline, live validation pending). Accidental exposure of held-out material in some engineers'
% repository-wide searches is recorded in their worklogs (F149, F151); open question whether to state it.

\subsection{Ablations of the harness}
\label{sec:setup-ablation}

The ablations test the prediction of Section~\ref{sec:method}, that removing one mechanism matters where nothing else
observes or bounds the consequence. Each variant removes one of the four maps and nothing else. \emph{No verification}
takes each effect from the skill's success flag and never looks at a carried item again. \emph{No facts} removes the
authoritative record from the context, leaving the model its own working plan and its last 12 steps. \emph{No guard}
checks no call. Its authority rule stays in the standing text, and an authorization request names the call for mission
control to decide. \emph{No context} replaces the rebuilt context with one accumulating conversation. With no variant
selected, all 1099 recorded steps of the 25 Experiment I runs of the proposed runtime re-render identically.
% F148 "Variants", render_check 1099/1099 over the 25 exp1b Ours runs; "last 12 steps": facts.py RECENT = 12.

Natural runs rarely exercise these mechanisms, so the ablations inject accidents into what the sensors and records
reported, never into simulator truth. J1 reports a carried item as gone from the jaws, at 17 points. J2 reports each of
the 8 failed presses as successful, and asks once right after it and once after a telemetry reading that shows the unit
unchanged. J3 restarts the reasoner after the first world-changing verified effect of each run, at 25 points. J4 starts
the L4 mission 755\,s later in the same pass, so Earth sets 170\,s after the start, at 3 points. Two probes without an
accident complete the design. One compares the full runtime with no context two thirds through each run (25 points).
The other compares it with no facts where an item's recorded location is older than the recent window (14 points).
% F148 "Accidents": J1 = first look at a carried item in each of 16 carrying runs + exp1's real drop = 17; J2 8 + 8;
% J3 25; J4 3 (755 s later, Earth-set 170 s after start); N 25; NF 14.

Stage one reads the Experiment I record for what each variant would have changed, with no model call. Stage two replays
each point offline. The exact system text, messages and palette of the point go to the Experiment I reasoning model three
times per arm. The decision is the first effect call, request to mission control or query of the robot's current state,
after applying up to six intervening plan calls offline. A rubric fixed before any model call classifies each decision,
and a reader blind to the arm scores those that need the rationale. Stage two comprised 684 replay chains and 828 model
calls. Stage three runs live from the recorded snapshot, only where stage two shows the variant clearly worse, with at
most 12 runs. J1 qualified for live runs of the full runtime and no verification at two instances. J3 runs live with no
context at two instances as a planned exception, since one decision cannot show a later repetition. Stage two sees only
the step after the accident.
% F148 "Rubric", "Amendment before the scored calls" (at most 6), "Stage 3" (at most 12); F149 (684 chains, 828 model
% calls, 240 decisions to the blind reader); F150 ("live: J1 (4) and J3 (2)"). Live runs pending the model gateway;
% results in Section results-ablation are TBD.
% Disagreement with research_content.md §5.7: it planned focused ablations on natural runs of a subset of Experiment I
% cases. F130 showed that natural runs would most likely give null results; the build uses injected accidents with
% offline counterfactual replay first and live runs only where offline shows a difference (user decision, F148).

\subsection{Metrics and statistics}
\label{sec:setup-metrics}

Autonomous mission success is the external evaluator's verdict with no artifact changed after the reveal. Human
independence follows ALFUS \citep{huang2007alfus}, computed identically for all methods from mission control's log and
the capability boundary's call windows. A mission spans from the first call's start to the later of the last call's end
and the last exchange with mission control. The human intervention ratio divides mission control's total decision time
by that span. Requests are counted by kind. The highest interaction level is the most human-directed ALFUS level at which
mission control answered. The re-authoring flag marks a new-procedure request, answered or queued. The idle fraction is
the share of the mission with a request open and no capability executing. The longest autonomous stretch is the longest
interval with no request open, the neglect time of \citet{olsen2003metrics}.
% evaluation/interaction.py measures() (mission span, HIR, idle_waiting_fraction, longest_autonomous_stretch_s,
% ground_reauthoring_required); ALFUS_ORDER.

The evaluator counts five classes of violation from simulator truth, by one rule for every method. An unintended contact
is a force above 10\,N on a robot body other than a wheel, outside a call whose contract ends in contact. A dropped item
comes to rest off every support, not held, more than 0.10\,m below where it last stood. A collateral change is another
unit's state changing, or an object that no call acts on moving by more than 1\,mm. The other classes are execution
during a latched safety fault and removal of a module from an energized unit without approval. A false completion is a
final report sent with every task of the working plan closed while the evaluator scores failure. The baselines send no
report, and a baseline plan that ends normally on a failed mission is stated in the results. Trace completeness counts
calls of the proposed runtime that carry proposal, guard outcome, execution result and verification. For PLEXIL it
counts boundary calls joined to their plan node by its identifier, and the checker does not cover
EUROPA\,$\rightarrow$\,PLEXIL. Experiment II adds active engineering time, the validation runs to the first pass, the
lines changed and held-out success before and after the extension.
% docs/evaluator_violations.md (10 N, SRB's own threshold; wheels excluded F109); DROP_HEIGHT_M = 0.10 (F143);
% DISPLACEMENT_NOISE_FLOOR_M = 1e-3; evaluation/rows.py false_completion (Experiment II reads the report texts,
% experiments/exp2/tables.py CLAIMS, F146). Trace: F133, evaluation/trace_completeness_check.py. Exp II columns:
% paper/tables/exp2_engineering.tex.

Success rates carry Wilson 95\% intervals \citep{wilson1927probable}, paired by instance. With samples this small the
comparison across coverage is descriptive, and no significance test was chosen after the data. Interaction measures are
means over the runs of each method and coverage. Metrics, failure taxonomy, exclusion rules and statistics were fixed
before the campaign. Each failed run gets one primary failure category, the first failure on its causal path. An
infrastructure failure, such as a simulator crash or a model-provider outage, is repeated and never scored. A physical
failure after the retry budget counts as a failed run.
% docs/preregistration.md §2 (eight categories), §3 (exclusion), §4 (Wilson, paired, descriptive). evaluation/tables.py
% interaction_table: means over runs. Disagreement: preregistration §4 says HIR, waiting and neglect time are reported as
% medians with bootstrap intervals; the generated tables report means without intervals (open question).
% Disagreement: research_content.md §5.5 metric 3 (human planning burden / preplanned transition fraction) is not
% computed (evaluation/README.md: pending); the re-authoring flag and the re-authoring study stand in for it.
% research_content.md §5.8 (real LunarBot demonstration) has no counterpart in the build yet.
